% P16-C4: structural and algebraic foundations for general identification.
\section{Latent-variable identification foundations}
\label{sec:p16-latent-identification-foundations}

The identification algorithm of Shpitser and Pearl uses kernels attached
to confounded components and reductions to ancestral subgraphs
\cite[Definitions~2--3 and Lemmas~4--8, pp.~1220--1224]{ShpitserPearl2006JointIdentification}.
This section supplies complete finite proofs of those foundations.
The structural model class, the information encoded by the observed
mixed graph, and the domains of conditional-ratio calculations are fixed
explicitly before those calculations are used.

\subsection{The independent-latent-input model class}

\begin{hbdefinition}{def:p16-latent-input-model}{Finite observed equations with independent latent roots}
Let \(N=\{1,\ldots,n\}\) index the observed variables \(V_i\),
with finite nonempty value sets \(S_i\). Their declared observed
parent sets \(\operatorname{pa}(i)\) form a directed acyclic graph.
Let \(\Lambda\) be a finite index set of latent inputs. Each
\(U_\alpha\) has a finite nonempty value set \(\mathcal U_\alpha\),
mass function \(r_\alpha\), and declared observed child set
\(C_\alpha\subseteq N\). The latent inputs are mutually independent:
\[
 r(u)=\prod_{\alpha\in\Lambda}r_\alpha(u_\alpha),
 \qquad u\in\mathcal U=\prod_{\alpha\in\Lambda}\mathcal U_\alpha.
\]
Put \(\Lambda_i=\{\alpha:i\in C_\alpha\}\). The total equations are
\[
 V_i=f_i(V_{\operatorname{pa}(i)},U_{\Lambda_i}),
 \qquad
 f_i:S_{\operatorname{pa}(i)}\times
          \prod_{\alpha\in\Lambda_i}\mathcal U_\alpha\longrightarrow S_i.
\]
An empty product contains one empty tuple. A listed argument may be
ignored by a particular function. Latent inputs are exogenous roots;
they have no equations using observed variables as parents.

The \emph{observed mixed graph} \(G\) has a directed edge
\(j\to i\) for each declared observed parent, and a bidirected
edge \(i\leftrightarrow j\) between distinct vertices sharing a
declared latent input. Directed acyclicity concerns the directed edges.
The bidirected edge records the common input. A latent input may have
one, two, or more children; its actual child set remains part of the
model data. One input shared by three children is not replaced here
by three independent pair-specific inputs.
For \(H\subseteq N\), the induced graph \(G[H]\) retains the
vertices \(H\) and both kinds of edge whose endpoints lie in \(H\).

This is the independent finite input class in the source's model
definition \cite[p.~1220]{ShpitserPearl2006JointIdentification}.
The source's observed set \(\mathbf V\) is indexed here by \(N\).
For a fixed mixed graph and observed value sets, the corresponding
model class permits all finite latent data and total functions with
those declared edges. Additional restrictions on latent child sets,
cardinalities, or distributions specify smaller model classes and
must be retained if they are imposed.
\end{hbdefinition}

\begin{hbdefinition}{def:p16-latent-q-kernel}{The intervention kernel Q[H]}
For \(H\subseteq N\) and a full assignment \(v\in S_N\), define
\[
 Q[H](v)=
 \mathbb P\bigl(V_H=v_H\mid
       \operatorname{do}(V_{N\setminus H}=v_{N\setminus H})\bigr).
\]
The argument \(v_H\) is the random output assignment and
\(v_{N\setminus H}\) is a fixed intervention parameter. The bracket
notation agrees with the source's \(Q[C]=P_{\mathbf v\setminus\mathbf c}(\mathbf c)\).
The source's complement notation there means the assigned values on
the observed complement of \(C\). One may also write the kernel
as \(Q[H](v_H;v_{N\setminus H})\) to separate output and parameter.
It does not mean conditioning on the observed outside assignment.
Write \(p(v)=Q[N](v)\). The empty kernel is \(Q[\varnothing]=1\).
Write \(p_H\) for the ordinary marginal of \(p\) on \(H\).
For disjoint observed sets \(X,Y\), the source notation \(P_x(y)\)
means \(\mathbb P(V_Y=y\mid\operatorname{do}(V_X=x))\), with
\(x\in S_X\) and \(y\in S_Y\).
\end{hbdefinition}

\begin{hbtheorem}{thm:p16-latent-kernel-semantics}{proved=annex:p16.latent-kernel-semantics}{A full-DAG representation and the exact kernel formula}
The model is represented by an acyclic structural model with the
latent inputs as explicit root vertices. For every \(H,v\),
\[
 Q[H](v)=\sum_{u\in\mathcal U}
       \left[\prod_{\alpha\in\Lambda}r_\alpha(u_\alpha)\right]
       \prod_{i\in H}
       \mathbf1_{\{v_i=f_i(v_{\operatorname{pa}(i)},u_{\Lambda_i})\}}.
\]
For each outside assignment this is a normalized mass function on
\(S_H\). It depends on outside values only at
\(\operatorname{pa}(H)\setminus H\), where
\(\operatorname{pa}(H)=\bigcup_{i\in H}\operatorname{pa}(i)\).
It is defined even if the outside assignment never occurs observationally.

For nonempty \(H\), fixing the outside equations and retaining only the observed
vertices in \(H\) gives a model of the same class. Its latent child
sets are \(C_\alpha\cap H\), its functions have the outside
parent arguments fixed to their assigned values, its observed mixed
graph is the induced graph \(G[H]\), and its distribution is
the kernel \(Q[H]\) at those parameters. This is a proved
restriction after intervention, not a general conversion of arbitrary
hidden endogenous graphs into this model class.
For empty \(H\), the retained observed assignment is the single
empty tuple and its kernel is one.
\end{hbtheorem}

\subsection{Districts and directed ancestry}

\begin{hbdefinition}{def:p16-latent-district}{Districts and the source C-components}
In the induced graph \(G[H]\), a \emph{bidirected path}
uses bidirected edges only. Two vertices of \(H\) are in the
same \emph{district} if they can be joined by such a path, including
the length-zero path from a vertex to itself. Write \(\mathcal D_H\)
for the resulting maximal classes; these are the source's maximal
\emph{C-components}. The empty graph has no districts, and an
isolated vertex is a singleton district.

Write \(\operatorname{an}_{G[H]}(A)\) for \(A\) together with
the vertices in \(H\) having a positive-length directed path to \(A\).
An \emph{ancestral subset within \(H\)} is \(A\subseteq H\)
that contains every directed ancestor in \(G[H]\) of its members.
Equivalently, every observed parent lying in \(H\) of a vertex
in \(A\) also lies in \(A\). Bidirected connectivity is not
used to define directed ancestry.
\index{district!causal graph}\index{C-component}\index{Q-factor}
\end{hbdefinition}

\begin{hblemma}{thm:p16-latent-district-partition}{proved=annex:p16.latent-district-partition}{The unique district partition and its latent input groups}
The districts uniquely partition \(H\). A nonempty induced subgraph
is bidirected-connected exactly when its bidirected edges contain a
spanning tree, giving the source's equivalent C-component definition.
For every latent input \(\alpha\), the nonempty set
\(C_\alpha\cap H\), when present, lies in exactly one district.
Thus the latent inputs with any child in \(H\) partition into
\[
 \Lambda_D=\{\alpha:\varnothing\ne C_\alpha\cap H\subseteq D\},
 \qquad D\in\mathcal D_H.
\]
The remaining inputs have no child in \(H\).
\end{hblemma}

\begin{hbtheorem}{thm:p16-latent-ancestral-reduction}{proved=annex:p16.latent-ancestral-reduction}{Ancestral summation of a Q kernel}
If \(A\) is ancestral within \(H\), then
\[
 \sum_{v_{H\setminus A}}Q[H](v)=Q[A](v).
\]
The left side holds \(v_{N\setminus H}\) and \(v_A\) fixed.
The right side is independent of the additionally displayed parameter
\(v_{H\setminus A}\), since these vertices are not parents of
\(A\). No positivity hypothesis is required.
\end{hbtheorem}

\begin{hbcorollary}{thm:p16-latent-target-reduction}{proved=annex:p16.latent-target-reduction}{Reducing a target to its remaining ancestors}
For disjoint observed sets \(X,Y\), put \(H=N\setminus X\)
and \(A=\operatorname{an}_{G[H]}(Y)\), with ancestry including
\(Y\). For every intervention assignment \(x\),
\[
 P_x(y)=\sum_{v_{A\setminus Y}}Q[A](v),
\]
where \(v_X=x\), and the values on \(H\setminus A\) can
be fixed arbitrarily. The right side is independent of those values.
Consequently fixing the additional coordinates \(H\setminus A\)
does not change this response law.

One may also first restrict the whole model to
\(A_0=\operatorname{an}_G(Y)\): its observed law is
the marginal \(p_{A_0}\), and the effect on \(Y\) requires only the
intervention coordinates \(X\cap A_0\). These are reductions
of specified models and kernels; they introduce no conditional ratio.
\end{hbcorollary}

\begin{hbtheorem}{thm:p16-latent-district-decomposition}{proved=annex:p16.latent-district-decomposition}{The district product for every intervention kernel}
For every \(H\subseteq N\) and full assignment \(v\),
\[
 Q[H](v)=\prod_{D\in\mathcal D_H}Q[D](v).
\]
Each factor is normalized on its own output set with the other
observed coordinates fixed. Directed parents can occur in other
districts, so this factorization is not an assertion of independence
of the observed district tuples. The identity remains true at zero
kernel values.
\end{hbtheorem}

\subsection{Recovering district factors from an available kernel}

\begin{hbdefinition}{def:p16-latent-prefixes}{Topological prefixes and positive conditional ratios}
For nonempty \(H\), choose an order \(h_1,\ldots,h_k\) in
which every directed parent within \(H\) precedes its child.
Let \(H_j=\{h_1,\ldots,h_j\}\), with \(H_0=\varnothing\),
and define
\[
 \rho_j(v)=\sum_{v_{H\setminus H_j}}Q[H](v),
 \qquad \rho_0=1.
\]
Outside values \(v_{N\setminus H}\) remain fixed. When
\(\rho_{j-1}(v)>0\), define
\(\theta_j(v)=\rho_j(v)/\rho_{j-1}(v)\).
This is a conditional mass in the intervened model with law \(Q[H]\).
For a district \(D\in\mathcal D_H\), put
\(m(D)=\max\{j:h_j\in D\}\).
\end{hbdefinition}

\begin{hblemma}{thm:p16-latent-prefix-chain}{proved=annex:p16.latent-prefix-chain}{Prefix kernels and the ordinary conditional chain}
Each \(H_j\) is ancestral within \(H\), and
\(\rho_j(v)=Q[H_j](v)\). The function \(\rho_j\) depends
only on the prefix assignment and outside parameters. On a positive
prefix, \(\theta_j\) sums to one over the next value. If
\(\rho_{k-1}(v)>0\), then
\[
 Q[H](v)=\prod_{j=1}^k\theta_j(v).
\]
The final numerator may be zero. No ratio with zero denominator is
defined by this construction.
\end{hblemma}

\begin{hbtheorem}{thm:p16-latent-district-recovery}{proved=annex:p16.latent-district-recovery}{The topological conditional-factor identity for a district}
For \(D\in\mathcal D_H\), at every assignment satisfying
\(\rho_{m(D)-1}(v)>0\),
\[
 Q[D](v)=\prod_{\{j:h_j\in D\}}
                \frac{\rho_j(v)}{\rho_{j-1}(v)}.
\]
All the displayed denominators are positive. In particular, if
\(Q[H]\) is strictly positive on \(S_H\) at the fixed outside
parameters, every district kernel is recovered there at all assignments.
For \(H=N\), a strictly positive observed law therefore identifies
all maximal district factors as
\[
 Q[D](v)=\prod_{\{j:h_j\in D\}}
             p(v_{h_j}\mid v_{H_{j-1}}).
\]
The more local positive-prefix condition also permits some zero
joint masses. A later recursive use must retain the available
parameter domain and justify every required denominator. The result
concerns maximal districts of the specified current graph; an
arbitrary smaller connected subset need not be identifiable.
\end{hbtheorem}

\subsection{Calculated kernels, reductions and limits}

\begin{hbdefinition}{def:p16-latent-foundation-example}{A three-variable model with two districts}
Name the observed vertices \(A,B,C\). Let independent latent inputs
\(U,E_A,E_B\) be bits with \(U\) fair,
\(\mathbb P(E_A=1)=1/5\), and \(\mathbb P(E_B=1)=1/4\).
Their declared child sets are \(\{A,C\}\), \(\{A\}\), and
\(\{B\}\), respectively. Set
\[
 A=U\oplus E_A,\qquad B=A\oplus E_B,\qquad
 C=\mathbf1_{\{B=1\text{ or }U=1\}}.
\]
The observed graph has \(A\to B\to C\) and \(A\leftrightarrow C\).
Write \(\alpha(0)=4/5,\alpha(1)=1/5\) and
\(\beta(0)=3/4,\beta(1)=1/4\).
The outcome \((u,e_A,e_B)\) has mass
\(\frac12\alpha(e_A)\beta(e_B)\).
\end{hbdefinition}

\begin{hbtheorem}{thm:p16-latent-foundation-calculation}{proved=annex:p16.latent-foundation-calculation}{The full district calculation, including zero numerators}
The full law and district kernels are
\[
\begin{array}{cc|cc}
 a&b&p(a,b,0)&p(a,b,1)\\\hline
 0&0&12/40&3/40\\
 0&1&0&5/40\\
 1&0&1/40&4/40\\
 1&1&0&15/40
\end{array}
\]
and
\[
 Q[\{B\}](a,b,c)=\beta(b\oplus a),\qquad
 Q[\{A,C\}](a,b,c)=
 \begin{cases}
 \frac12\alpha(a\oplus c),&b=0,\\
 \frac12\mathbf1_{\{c=1\}},&b=1.
 \end{cases}
\]
The maximal districts are \(\{A,C\}\) and \(\{B\}\), and
the product of these kernels gives every cell above. In order
\(A,B,C\), the prefix masses are
\(\rho_1=1/2\) and \(\rho_2=\frac12\beta(b\oplus a)>0\).
Consequently the district recovery formula holds at every assignment,
including the zero cells with \(b=1,c=0\).

In \(H=\{A,C\}\), both singletons are directed-ancestral, though
the vertices remain bidirected-connected. Summation gives
\[
 Q[\{C\}](a,b,1)=
 \begin{cases}1/2,&b=0,\\1,&b=1,\end{cases}
 \qquad Q[\{A\}]=1/2.
\]
This exhibits external kernel parameters, ancestral reduction, and
district extraction in the same finite model.
\end{hbtheorem}

\begin{hbtheorem}{thm:p16-latent-foundation-limits}{proved=annex:p16.latent-foundation-limits}{Support and latent-incidence restrictions remain substantive}
The following concrete limits hold.
\begin{enumerate}
\item With allowed graph \(A\to B\), the models \(A=0,B=0\)
      and \(A=0,B=A\) have the same observed point mass, but
      their kernels for \(B\) at \(A=1\) differ. The required
      prefix denominator is zero.
\item With graph \(A\to Y\) and \(A\leftrightarrow Y\),
      let independent fair roots \(U,T\) have child sets
      \(\{A,Y\}\) and \(\{Y\}\). The two models
      \(A=U,Y=\mathbf1_{\{T=1\text{ or }U\oplus A=1\}}\)
      and \(A=U,Y=T\) have the same strictly positive observed
      law, but \(Q[\{Y\}](Y=1;A=a)\) equals \(3/4\) in
      the first and \(1/2\) in the second. The singleton is
      connected, but is not a maximal district of the full graph.
\item A fair root shared by three observed children with equations
      \(A=B=C=U\) gives probability \(1/2\) each to
      \(000,111\). Three independent fair pair-specific roots,
      with each observed variable the exclusive-or of its two incident
      roots, instead give probability \(1/4\) each to
      \(000,011,101,110\). Both declared mixed graphs are the
      bidirected triangle. That particular root replacement does
      not preserve the law.
\end{enumerate}
The third example compares specified constructions; it asserts no
general impossibility theorem about every other pairwise representation.
\end{hbtheorem}

The complete identification algorithm and its recursive invariants are
supplied in \cref{thm:p16-id-invariants,thm:p16-id-termination,
thm:p16-id-success-soundness}. The hedge construction and its proved
completeness classes follow in
\cref{thm:p16-hedge-failure-extraction,thm:p16-hedge-completeness}.
Every invocation retains its stated value sets and support conditions.
